\documentclass[14pt,a4paper,oneside]{extbook}
\usepackage[utf8]{vietnam}
\usepackage{amsmath,amssymb,mathrsfs,fancyhdr,enumerate,multirow,makecell}
\usepackage{tkz-euclide,tikz-3dplot,tikz,tkz-tab}
\usetikzlibrary{arrows.meta,calc,intersections,angles,quotes,backgrounds,shapes.geometric,patterns,through,shadows,positioning}
\usepackage{yhmath}
\usepackage[top=0.7cm, bottom=0.7cm, left=1.0cm, right=1.0cm, headheight=16pt, headsep=3mm, footskip=5mm]{geometry}
\usepackage[hidelinks,unicode]{hyperref}
\usepackage{qrcode}
\usepackage[many]{tcolorbox}
\usepackage{fontawesome}

\pagestyle{fancy}
\fancyhf{}
\fancyhead[L]{\footnotesize\bfseries\color{teal!85!black}\faBook~SỔ TAY KIẾN THỨC TOÁN 5}
\fancyhead[R]{\footnotesize\bfseries\color{blue!85!black}Lớp Toán Thầy Lĩnh -- loptoanthaylinh.vn}
\fancyfoot[C]{\footnotesize\bfseries\color{gray!80!black}-- Sổ tay kiến thức: Bài 6. Cộng, trừ hai phân số --}
\renewcommand{\headrulewidth}{0.4pt}

\begin{document}

% --- BANNER TIÊU ĐỀ BÀI HỌC ---
\noindent
\begin{tcolorbox}[
  enhanced,
  colback=teal!8!white,
  colframe=teal!80!black,
  arc=3mm,
  boxrule=1.1pt,
  width=\linewidth,
  halign=center,
  top=2.2mm, bottom=2.2mm,
  left=3mm, right=3mm,
  drop fuzzy shadow
]
  {\large\bfseries\color{teal!95!black}\faBookmark\ BÀI 06. CỘNG, TRỪ HAI PHÂN SỐ}\\[2pt]
  {\scriptsize\color{black!75}\itshape Kỹ thuật chọn mẫu số chung nhỏ nhất, cộng trừ với số tự nhiên trên 1 trang A4}
\end{tcolorbox}

\vspace{3.2mm}
% --- KHỐI 1: QUY TẮC CƠ BẢN ---
\noindent
\begin{minipage}[t]{0.55\textwidth}
  \begin{tcolorbox}[
    enhanced,
    title={\footnotesize\bfseries\color{white}\faKey\ 1. Quy trình cộng, trừ 2 phân số khác mẫu},
    colback=blue!2!white,
    colframe=blue!75!black,
    arc=2.5mm, boxrule=0.8pt, top=1.8mm, bottom=1.8mm, left=2.2mm, right=2.2mm
  ]
    \footnotesize
    \textbf{Bước 1:} Tìm \textbf{Mẫu số chung} (MSC) thích hợp.\\
    $\triangleright$ Tốt nhất chọn số nhỏ nhất chia hết cho cả hai mẫu.\\
    \textbf{Bước 2:} Tìm thừa số phụ của từng mẫu số.\\
    \textbf{Bước 3:} Quy đồng tử số tương ứng.\\
    \textbf{Bước 4:} Thực hiện phép cộng (hoặc trừ) tử số và giữ nguyên MSC.\\
    \textbf{Bước 5:} Rút gọn kết quả về tối giản (nếu chưa tối giản).
  \end{tcolorbox}
\end{minipage}%
\hfill
\begin{minipage}[t]{0.43\textwidth}
  \begin{tcolorbox}[
    enhanced,
    title={\footnotesize\bfseries\color{white}\faSliders\ 2. Cộng, trừ với số tự nhiên},
    colback=teal!2!white,
    colframe=teal!75!black,
    arc=2.5mm, boxrule=0.8pt, top=1.8mm, bottom=1.8mm, left=2mm, right=2mm
  ]
    \footnotesize
    Viết số tự nhiên dưới dạng phân số có mẫu số bằng MSC:\\
    $\bullet$ $1 + \dfrac{3}{5} = \dfrac{5}{5} + \dfrac{3}{5} = \dfrac{8}{5}$.\\[2pt]
    $\bullet$ $2 - \dfrac{3}{4} = \dfrac{8}{4} - \dfrac{3}{4} = \dfrac{5}{4}$.\\[2pt]
    $\bullet$ $\dfrac{11}{3} - 3 = \dfrac{11}{3} - \dfrac{9}{3} = \dfrac{2}{3}$.\\[2pt]
    \textit{Mẹo nhanh:} $a \pm \dfrac{b}{c} = \dfrac{a \times c \pm b}{c}$.
  \end{tcolorbox}
\end{minipage}

\vspace{3mm}

% --- KHỐI 2: KỸ THUẬT CHỌN MẪU SỐ CHUNG TỐI ƯU ---
\noindent
\begin{tcolorbox}[
  enhanced,
  title={\footnotesize\bfseries\color{white}\faLightbulbO\ 3. Ba trường hợp tìm Mẫu số chung (MSC) thông minh},
  colback=purple!2!white,
  colframe=purple!70!black,
  arc=2.5mm, boxrule=0.8pt, top=1.8mm, bottom=1.8mm, left=2.5mm, right=2.5mm
]
  \footnotesize
  \begin{tabular}{ll}
    $\bullet$ \textbf{Trường hợp 1 (Mẫu lớn chia hết cho mẫu bé):} Lấy luôn mẫu lớn làm MSC.\\
    \quad \textit{VD:} $\dfrac{5}{6} + \dfrac{1}{18}$: Vì $18$ chia hết cho $6$ nên chọn $\text{MSC} = 18 \rightarrow \dfrac{5 \times 3}{18} + \dfrac{1}{18} = \dfrac{16}{18} = \dfrac{8}{9}$.\\
    $\bullet$ \textbf{Trường hợp 2 (Hai mẫu không có ước chung $>1$):} Lấy tích hai mẫu làm MSC.\\
    \quad \textit{VD:} $\dfrac{2}{3} + \dfrac{1}{5}$: Chọn $\text{MSC} = 3 \times 5 = 15 \rightarrow \dfrac{10}{15} + \dfrac{3}{15} = \dfrac{13}{15}$.\\
    $\bullet$ \textbf{Trường hợp 3 (Cùng chia hết cho một số):} Nhân mẫu lớn với $2, 3\dots$ xem có chia hết cho mẫu kia không.\\
    \quad \textit{VD:} Mẫu $8$ và $12$: Thử $12 \times 2 = 24$ (chia hết cho $8$) $\rightarrow$ Chọn $\text{MSC} = 24$ (không cần lấy $8 \times 12 = 96$).
  \end{tabular}
\end{tcolorbox}

\vspace{3mm}

% --- KHỐI 3: SƠ ĐỒ TƯ DUY ---
\noindent
\begin{tcolorbox}[
  enhanced,
  title={\footnotesize\bfseries\color{white}\faSitemap\ Sơ đồ tư duy: Các bước cộng, trừ phân số},
  colback=yellow!2!white,
  colframe=teal!85!black,
  arc=3mm, boxrule=0.9pt, top=1.8mm, bottom=1.8mm, left=1.5mm, right=1.5mm,
  halign=center
]
  \begin{tikzpicture}[
    scale=0.92,
    every node/.style={transform shape},
    s1/.style={rectangle, rounded corners=4pt, fill=blue!15, draw=blue!70!black, text=blue!90!black, font=\bfseries\scriptsize, inner sep=4pt, align=center},
    s2/.style={rectangle, rounded corners=4pt, fill=teal!15, draw=teal!70!black, text=teal!90!black, font=\bfseries\scriptsize, inner sep=4pt, align=center},
    s3/.style={rectangle, rounded corners=4pt, fill=orange!15, draw=orange!70!black, text=orange!90!black, font=\bfseries\scriptsize, inner sep=4pt, align=center},
    s4/.style={rectangle, rounded corners=4pt, fill=green!15, draw=green!70!black, text=green!90!black, font=\bfseries\scriptsize, inner sep=4pt, align=center}
  ]
    \node[s1] (s1) at (-6.3, 0) {1. Quan sát mẫu số\\(Cùng hay Khác mẫu?)};
    \node[s2] (s2) at (-2.1, 0) {2. Khác mẫu:\\Quy đồng mẫu chung};
    \node[s3] (s3) at (2.1, 0) {3. Cộng / trừ tử số\\Giữ nguyên mẫu};
    \node[s4] (s4) at (6.3, 0) {4. Rút gọn kết quả\\(Về PSTG)};

    \draw[teal!85!black, thick, -{Stealth[length=2.5mm]}] (s1) -- (s2);
    \draw[teal!85!black, thick, -{Stealth[length=2.5mm]}] (s2) -- (s3);
    \draw[teal!85!black, thick, -{Stealth[length=2.5mm]}] (s3) -- (s4);
  \end{tikzpicture}
\end{tcolorbox}

\vspace{3mm}

% --- KHỐI 4: TÍNH NHANH & BẪY TRẮC NGHIỆM ---
\noindent
\begin{minipage}[t]{0.54\textwidth}
  \begin{tcolorbox}[
    enhanced,
    title={\footnotesize\bfseries\color{white}\faBolt\ Tính chất phân phối và nhóm số hạng},
    colback=green!3!white,
    colframe=green!60!black,
    arc=2.5mm, boxrule=0.8pt, top=1.8mm, bottom=1.8mm, left=2.2mm, right=2.2mm
  ]
    \footnotesize
    Nhóm các phân số có cùng mẫu lại với nhau:\\
    \textit{VD:} $\left(\dfrac{3}{8} + \dfrac{2}{5}\right) + \dfrac{5}{8} = \left(\dfrac{3}{8} + \dfrac{5}{8}\right) + \dfrac{2}{5} = 1 + \dfrac{2}{5} = \dfrac{7}{5}$.\\
    Tránh quy đồng cả 3 phân số cùng một lúc!
  \end{tcolorbox}
\end{minipage}%
\hfill
\begin{minipage}[t]{0.44\textwidth}
  \begin{tcolorbox}[
    enhanced,
    title={\footnotesize\bfseries\color{white}\faExclamationTriangle\ Mẹo tránh bẫy trắc nghiệm},
    colback=red!2!white,
    colframe=red!70!black,
    arc=2.5mm, boxrule=0.8pt, top=1.8mm, bottom=1.8mm, left=2.2mm, right=2.2mm
  ]
    \footnotesize
    $\bullet$ Luôn kiểm tra xem phân số ban đầu đã \textbf{tối giản} chưa. Nếu chưa, hãy rút gọn trước để việc quy đồng mẫu trở nên cực kì nhẹ nhàng.\\
    $\bullet$ Quên rút gọn đáp số là lỗi bị trừ điểm phổ biến nhất!
  \end{tcolorbox}
\end{minipage}

\end{document}
